\section{La decisión histórica de CUDA: convertir la capacidad de paralelismo de un «truco gráfico» en «infraestructura de uso general»}

\subsection{De «disfrazar un problema científico como un problema gráfico» a una interfaz estandarizada}

Jensen recordó lo incómodo que era el periodo inicial de GPGPU (General-purpose computing on GPU, computación de propósito general en GPU): los investigadores tenían que «engañar» al flujo gráfico para que la GPU creyera que estaba haciendo una tarea de gráficos, y así aprovechar su capacidad de procesamiento paralelo. Esta vía permitía hacer funcionar experimentos, pero no podía dar lugar a un ecosistema de ingeniería.

El valor decisivo de CUDA está en esto: desacopló la programación de GPU de la semántica de las API gráficas y la convirtió en un lenguaje general y una cadena de herramientas de desarrollo de uso general. En la entrevista él subraya que «puedes usar C para decirle a la GPU qué hacer», lo cual, en esencia, es elevar la «función de producción del desarrollador».

\begin{figure}[H]
\centering
\includegraphics[width=0.88\textwidth]{frame-cuda.jpg}
\caption{CUDA como capa de software que abre al exterior el Program Execution de la GPU \protect\footnotemark}
\end{figure}
\footnotetext{Intervalo de tiempo en el video: 00:08:29--00:08:40.}

\begin{importantbox}{If you do not build it they cannot come}
Jensen da la lógica central de la inversión en plataformas: \textit{``If you build it, they might not come; but if you don't build it, they can't come.''}

Esta es una inversión típica en una plataforma anticipada: primero se soporta la incertidumbre y después se obtiene el techo del ecosistema. Cuando CUDA nació no tenía un «contrato asegurado con un gran cliente», pero sentó por adelantado la base de software que después necesitaría el auge de la IA.
\end{importantbox}

\subsection{Por qué NVIDIA se atrevió a apostar tan fuerte: el mercado de producción masiva ofrecía «viabilidad económica»}

Un punto que la entrevista suele pasar por alto, pero que es crucial, es que Jensen no era un «idealista ciego». Mencionó reiteradamente «el tamaño del mercado de los videojuegos», porque eso significaba que la GPU se convertiría en un procesador paralelo de alta producción a escala mundial. Una alta producción no trae consigo solo resonancia de mercadotecnia, sino dos restricciones firmes:
\begin{itemize}
    \item La curva de costos puede bajar lo suficiente para que instituciones de investigación y equipos emprendedores puedan costearla.
    \item El ritmo de iteración puede mantenerse, sin que el ecosistema de software sufra una ruptura generacional.
\end{itemize}

\begin{warningbox}{Error común: el valor de CUDA está solo en la «mejora de rendimiento»}
Si a CUDA se le entiende solamente como «hacer que un programa corra un poco más rápido», se subestima su significado histórico. El verdadero valor de CUDA es haber convertido la computación paralela de un «trabajo manual para expertos» en una «capacidad industrial que se puede enseñar, reutilizar y escalar». Este tipo de transferencia de capacidad suele generar, con el tiempo, más interés compuesto que una sola mejora de rendimiento.
\end{warningbox}

\subsection{Resumen de la sección}
CUDA no fue el lanzamiento de una herramienta, sino un compromiso organizacional a nivel de plataforma. Redefinió la GPU: de «chip gráfico» pasó a ser la «base de computación paralela de uso general», y con ello proporcionó una condición previa para la posterior industrialización de la IA.
